%=============================================================================
% Ambient Sense: Multi-Band Passive RF Sensing on Commodity Smartphones
% IEEE two-column manuscript — LaTeX source (revised for flow, consistency,
% and length).
% Compile with:  pdflatex Ambient_Sense_Paper.tex  (run twice for refs)
% Requires: the IEEEtran document class and the 'cite' package.
%=============================================================================
\documentclass[journal,10pt]{IEEEtran}

\usepackage{cite}          % numeric citations with range compression
\usepackage{graphicx}      % \includegraphics
\usepackage{amsmath,amssymb}
\usepackage{textcomp}
\usepackage{array}
\usepackage{booktabs}
\usepackage{url}
\usepackage[hidelinks]{hyperref}

\graphicspath{{figures/}}  % images live in the figures/ subfolder

\begin{document}

%---------------------------------------------------------------------------
% Title block
%---------------------------------------------------------------------------
\title{Ambient Sense: Multi-Band Passive RF Sensing on Commodity Smartphones}

\author{%
\IEEEauthorblockN{Prepared for submission to IEEE --- full manuscript draft}
\IEEEauthorblockA{Device-Free dBm-Based Ranging, Scan-Epoch Change Detection,
Two-Node Transit Counting, Mesh-Accrued Accuracy and Gamified Community
Reporting}
}

\markboth{Journal of \LaTeX\ Class Files,~Vol.~XX, No.~X, 2026}%
{Author \MakeLowercase{\textit{et al.}}: Ambient Sense}

\maketitle

%---------------------------------------------------------------------------
% Abstract & Index Terms
%---------------------------------------------------------------------------
\begin{abstract}
Device-free passive sensing typically requires dedicated radar arrays,
synchronised infrastructure, or camera feeds that raise privacy concerns. We
present Ambient Sense, a zero-infrastructure sensing framework that
repurposes the radios already inside a commodity smartphone --- BLE and
dual-band Wi-Fi --- into a passive, multi-band RF sentinel. By fusing
per-link attenuation dynamics across three carrier bands, Ambient Sense
detects and classifies non-transmitting pedestrians and vehicles without
target cooperation. To suppress environmental fluctuation and quiet-link
false alarms, the system compares timed field snapshots (``scan epochs'')
against stable, median-filtered baselines rather than raw samples, and gates
distance estimation at an explicit noise floor so that ambiguous links
publish a range \emph{band} instead of an unreliable point estimate.
Independent phones can further federate into a cooperative mesh, where
crossing events confirmed by multiple nodes extend the effective sensing
radius and a gamification layer rewards devices whose measurements agree
with the group. Every constant reported here was obtained by executing a
model-based Monte-Carlo evaluation harness that replays synthetic scenes
through the production signal-processing pipeline and rejects any candidate
configuration that does not win on complete scenes. Across this harness, the
system identifies crossing entities to $\approx$12~m and detects presence to
$\approx$16~m from a single phone, extends detection toward 20--40~m through
mesh cooperation, and reports --- rather than hides --- the ranges at which
the underlying physics stops supporting a confident estimate. As all results
here are model-based, on-device field validation is identified as the
primary next step.
\end{abstract}

\begin{IEEEkeywords}
device-free sensing, passive RF sensing, received signal strength indication
(RSSI), Bluetooth Low Energy, Wi-Fi sensing, channel state information,
device-free localization, people/flow counting, change detection, mobile
crowdsensing, gaming incentives.
\end{IEEEkeywords}

%---------------------------------------------------------------------------
% I. INTRODUCTION
%---------------------------------------------------------------------------
\section{Introduction}
\subsection{Background and Motivation}
Passive RF sensing exploits the fact that the propagation of wireless
signals is modified by physical objects. When a human body or a vehicle
crosses a radio link, it attenuates, scatters and diffracts the
electromagnetic field, and the received signal strength dips and recovers
with a characteristic signature. Because the observer requires no device
attached to the target, this paradigm is called device-free sensing, and it
underpins applications ranging from indoor localization and
ambient-assisted living to occupancy analytics \cite{wilson2010}. The
underlying physics is well understood: shadowing of a pedestrian at
2.4~GHz can reach several decibels with a lateral profile on the order of a
wavelength, and both the depth and the temporal texture of the attenuation
encode information about the entity that caused it~\cite{januszkiewicz2018}.
Modern systems build on statistical inference over RSSI or
channel-state-information (CSI) sequences, spanning radio tomographic
imaging (RTI) and its learned extensions~\cite{alippi2016,li2024,wu2024,cnnrti2025}
as well as activity-recognition pipelines that classify the motion
itself~\cite{wong2024,sanchez2020}.

A second observation is that smartphones already contain the radios such
sensing requires. A modern handset carries a BLE transceiver and a dual-band
(2.4~GHz and 5~GHz) Wi-Fi interface, and can observe the received signal
strength of every advertising beacon and access point (AP) in range
\emph{without} connecting to it. This converts the phone from a
communication terminal into a multi-band passive sensor, where each observed
link (phone--beacon or phone--AP) becomes a measurement ray whose
attenuation state can be monitored in real time. Beacon-level RSSI from
unmodified receivers has been shown to be usable as a sensing
signal~\cite{ramirez2021,milano2024,bai2020}, relocating the sensing problem
from dedicated node deployments~\cite{wilson2010,alippi2016} to software.

A third observation is economic: dedicated sensing hardware is costly, and
camera-based alternatives raise privacy concerns~\cite{khan2023,zhou2020,dmip2023}.
Smartphone-based passive sensing can be deployed at near-zero marginal cost
by people who already carry the device, and when several phones in an area
cooperate --- an \emph{Ambient Sense mesh} --- the geometric diversity of
their measurement rays can increase both accuracy and spatial
reach~\cite{li2024}. Mobile-crowdsensing research shows that such
cooperation must be engineered: participants contribute when mechanisms are
transparent, rewarded, and socially engaging~\cite{zhang2022,wang2018}.
Gamification is therefore treated here as incentive infrastructure rather
than decoration --- at the cost of making the \emph{trustworthiness} of
contributed measurements an explicit engineering concern.

\subsection{Problem Statement}
The gap addressed here is between what passive RF sensing demonstrably
achieves in the laboratory and what a single, unmodified smartphone held by
an ordinary user can deliver in the field. Three difficulties compound.
First, detection is physics-limited: occlusion attenuation decays with
distance while environmental fluctuation grows with range, so beyond some
distance the event becomes statistically indistinguishable from noise; a
system that does not quantify this limit misleads its user. Second, RSSI
range estimation is notoriously ambiguous --- the log-distance model that
converts decibels into metres is, on weak links, extremely ill-conditioned,
and an unconstrained estimate can be wrong by tens of
metres~\cite{ramirez2021,milano2024}. Third, cooperative mode introduces an
integrity problem: a device can report forged beacons or forged
measurements, and the mesh mean --- the very quantity that gamified
reliability rewards --- becomes an attack
surface~\cite{casar2022,antonioli2020}. Prior work frequently reports
accuracy without stating the distance at which the signal-to-noise ratio
actually supports it, conflates classification with detection, or assumes
infrastructure (synchronised nodes, CSI-capable chipsets) that phones do not
have.

\subsection{Approach and Contribution}
This paper presents \textbf{Ambient Sense}, a deployable native smartphone
application that performs passive sensing using only the radios already
present in the phone. Its behaviour is validated by a model-based Monte
Carlo harness that replays synthetic scenes through the production
signal-processing code; every constant reported in Section~III --- classifier
weights, trigger floors, the ranging gate, the epoch comparison rule, and the
mesh tolerance --- was accepted only after winning an A/B test on
\emph{complete scenes}, not on row-level aggregates. This falsification
protocol lets the paper state, with equal precision, both the configuration
that shipped and the configurations the evidence rejected (Section~III.D).

The system contributes four components, detailed fully in Section~II:
\begin{itemize}
  \item A \textbf{multivariate, multi-band fusion classifier} that combines
  per-link dip depth, duration, texture, variance, prominence and multi-link
  coherence across BLE, Wi-Fi~2.4~GHz and Wi-Fi~5~GHz into a
  pedestrian/vehicle decision, with per-band accuracy reported separately
  rather than pooled.
  \item A \textbf{gated dBm-to-metre ranging} component that inverts the
  log-distance model but refuses a point estimate below an explicit noise
  floor, publishing a radial \emph{band} instead of a misleadingly precise
  number on weak links.
  \item \textbf{Scan-epoch change detection} that compares two timed
  snapshots of the field using a stable, median-filtered baseline rather than
  a raw sample-to-sample delta, which is what suppresses false alarms.
  \item A \textbf{two-node transit counter} and a \textbf{cooperative mesh}
  in which multiple phones confirm each other's crossings, extending
  effective range at the cost of requiring multi-node agreement.
\end{itemize}
The paper's second contribution is methodological: it reports not only what
Ambient Sense achieves, but the specific distances, signal levels and
configurations at which it deliberately declines to estimate.

\subsection{Related Work}
\textbf{Device-free sensing and radio tomography.} RTI reconstructs an
attenuation field from link measurements over a node
network~\cite{wilson2010}; later work removes the need for training data,
recovers attenuation images generatively, and extends to multi-target and
outdoor scenes~\cite{alippi2016,wu2024}. RTI has since been integrated with
reconfigurable intelligent surfaces to roughly halve the required node
density~\cite{li2024}, and transformer/latent-variable successors improve
robustness under composite noise~\cite{wu2024}. CNN-based reconstruction from
Wi-Fi RSSI on ESP32 meshes reports over 92\,\% localization accuracy in
controlled geometry~\cite{cnnrti2025} --- a useful upper bound for what
RTI-style inversion can reach with known node geometry, and a reminder that
phone-held sensing lacks exactly that controlled geometry.

\textbf{Activity recognition from Wi-Fi.} CSI-based classifiers and recent
high-dimensional, multi-transceiver datasets have widened the range of
recognisable actions~\cite{wong2024}, and fusing CSI amplitude with RSSI from
a single dual-band AP has achieved sub-metre indoor
localization~\cite{sanchez2020}. These systems typically assume CSI
extraction or controlled capture hardware and target \emph{deliberately
performed} activities; Ambient Sense instead targets opportunistic detection
of people passing by, where only link-level RSSI is observable.

\textbf{Counting and flow estimation.} Device-free counting has been
demonstrated with CSI and deep networks~\cite{khan2023}, Doppler-based flow
estimation~\cite{zhou2020}, and commodity Wi-Fi counting/identification of
moving persons~\cite{dmip2023}. The common trajectory in this strand is the
move from counting toward flow; the two-node transit counter here realises
the same barrier-coverage geometry on commodity Bluetooth, without dedicated
infrastructure.

\textbf{Ranging and shadowing physics.} RSSI ranging requires integrity
monitoring against unstable beacons and missing
links~\cite{ramirez2021,bai2020}, and human-body shadowing at 2.4~GHz has
been characterised and shown to be compensable in wearable and
sensor-network settings~\cite{januszkiewicz2018}. Fisher-information analysis
of device-free localization in integrated sensing-and-communication networks
shows that positional information along a link collapses beyond a certain
SNR~\cite{li2024}, which supports the decision below to report distance
\emph{bands} rather than points at long range.

\textbf{Incentives and data integrity.} Quality-weighted and reputation-based
incentive design outperforms raw-volume rewards in mobile
crowdsensing~\cite{zhang2022,wang2018}. On the adversary side, BLE security
has been surveyed comprehensively~\cite{casar2022}, downgrade attacks on key
negotiation are known~\cite{antonioli2020}, and physical-layer fingerprints
rather than identifiers can support spoofing
detection~\cite{cai2024,mobisys2024,hensler2019}. These results bound what a
future adversarial-resilience module for the Ambient Sense mesh would need
to defend against; the present mesh design (Section~II.H) provides
majority-agreement weighting but does not yet implement such a defence, and
we return to this in Section~IV.

\subsection{Paper Organization}
Section~II presents the system design: the sensing model, the DSP chain,
multivariate fusion, ranging and gating, scan epochs, mesh agreement, the
two-node counter, gamification, and the privacy boundary. Section~III
reports the evaluation harness and the per-radius accuracy, epoch-detection,
ranging, mesh, and transit results drawn from it. Section~IV concludes and
outlines future work, including the on-device validation and adversarial
hardening that the current results do not yet cover.

%---------------------------------------------------------------------------
% II. SYSTEM DESIGN
%---------------------------------------------------------------------------
\section{System Design and Architecture}

\textbf{A. Sensing model and bands.} Ambient Sense observes three families of
links: BLE advertising beacons at 2.4~GHz, Wi-Fi APs at 2.4~GHz, and Wi-Fi
APs at 5~GHz. Received power on each link is modelled by log-distance path
loss, $\mathrm{RSSI}(d) = \mathrm{PL}_0 - 10n\log_{10}(d)$, with per-band
reference parameters --- path-loss exponent $n$ and 1~m reference strength
$R_{1\,\mathrm{m}}$ --- set to $(n, R_{1\,\mathrm{m}}) = (2.4, -42)$~dBm for
BLE, $(2.6, -38)$~dBm for Wi-Fi~2.4~GHz, and $(2.9, -35)$~dBm for
Wi-Fi~5~GHz. Environmental fluctuation is modelled as lognormal scintillation
whose standard deviation grows with range (0.45--0.70~dB at short range,
increasing by $\approx$0.055~dB/m). An occlusion is a Gaussian-profiled
attenuation $A_{\mathrm{eff}}\exp(-\mathrm{lat}^2/2\sigma^2)$ whose amplitude
decays with closest-approach distance through a Fresnel-efficiency term
$(1+d/10)^{-0.6}$, with lateral width $\sigma = 0.35$~m for a pedestrian and
$1.15$~m for a vehicle; vehicles scale the occlusion to $\approx$25~dB
(capped at 34~dB) and modulate it with a 5.5~Hz flutter, while pedestrians
carry a 3.1~Hz body wobble of $\approx$0.9~dB. This physics forms the ground
truth of the evaluation harness and mirrors, at first order, measured human
shadowing at 2.4~GHz~\cite{januszkiewicz2018}.

\begin{figure}[htb!]
  \centering
  \includegraphics[width=8.7cm]{figures/fig_modern_01_architecture.jpg}
  \caption{Functional decomposition of the Ambient Sense sensing pipeline.
  Each stage produces the artefacts (per-link dBm tables, epoch tables,
  transit counts and mesh ledgers) that the evaluation harness replays.}
  \label{fig:method}
\end{figure}

\textbf{B. Signal-processing chain.} A \texttt{RingWindow} keeps recent
samples; \texttt{Stats} maintains the median and robust dispersion;
\texttt{BaselineTracker} estimates the adaptive no-event level of each link;
and \texttt{DipDetector}/\texttt{PeakDetector} segment the stream into
events. Three detector-level rules emerged from measured failures: a
\emph{noise-aware trigger floor} (a dip must exceed a floor that scales with
the link's own measured fluctuation, so a quiet link triggers on a smaller
dip than a noisy one), a \emph{sustained-trigger gate} (the trigger must
persist for a minimum number of samples, removing isolated glitches), and the
\emph{stable-level policy} described in item~F below. Fig.~\ref{fig:method}
shows the resulting decomposition.

\textbf{C. Features.} For every detected event on every link, the engine
computes dip \emph{depth} (dB), \emph{duration} (ms), \emph{texture} (the
jitter ratio distinguishing a vehicle's smooth block from a pedestrian's
wobble), \emph{variance}, \emph{prominence} (depth relative to competing
peaks in the window), and \emph{coherence} (the number of links registering
the event inside a common interval --- the multi-link agreement that later
becomes the mesh primitive). Pedestrians produce 5--12~dB dips with jittery
texture; vehicles produce 15--26~dB dips with smooth, long blockages; and
confusion concentrates in the 4--13~dB band.

\textbf{D. Multivariate multi-band fusion.} \texttt{EntityClassifier}
applies a learned weight vector (shipped version v1.1) over each band's
features separately and fuses the bands at the crossing level. BLE is the
fast, dense human signature, sampling closer to the body and more
frequently; Wi-Fi~5~GHz contributes the strongest per-link attenuation
($A_0 = 11.5$~dB at reference range) over longer links; Wi-Fi~2.4~GHz
contributes intermediate-distance coverage with the lowest fluctuation
floor. The fusion is deliberately asymmetric: BLE carries event typing while
the Wi-Fi bands extend detection to 58--68~m link legs that BLE alone cannot
reach (Section~III.A quantifies the resulting per-band accuracy split). The
fused output is an entity type (pedestrian/vehicle) with a confidence and an
event record.

\textbf{E. dBm-based ranging and the quality ladder.} Ranging inverts the
log-distance model under an explicit honesty contract. Quantised signal
levels form a quality ladder from strong links ($-20$~dBm) down to a
$-90$~dBm noise floor~\cite{ramirez2021,milano2024,bai2020}. Estimation is
gated at $-75$~dBm: below this, the inverse function is so ill-conditioned
that an unconstrained point estimate can be off by 14--23~m (Section~III.B),
so the system refuses a point and instead publishes a radial band
$f \in [0.15, 0.92]$ of the raw distance. Links at or below $-90$~dBm are
excluded from range and accuracy computations entirely and retained only for
multi-link agreement.

\textbf{F. Scan-epoch change detection.} Rather than comparing successive
samples, the engine treats the field as two timed snapshots on a
$\approx$10~s moving-object time scale. Each epoch stores a timed record
(timestamp, place, geo, per-link dBm table); the engine then resets its
baselines and ring windows, rescans, and computes the difference from
\emph{stable levels} --- the median of the last 20 samples --- or a
transient dip of $\geq 3$~dB. This rule is deliberate: single-sample
comparisons have $\approx$1~dB standard deviation and produced a 25--62\,\%
false-positive rate in testing, whereas the 20-sample median collapses that
fluctuation before any comparison is made (Section~III.B).

\textbf{G. Two-node transit counter.} Two phones form a barrier across a
passage or road using a Classic Bluetooth RFCOMM link and a BLE L2CAP link.
The counter uses a hybrid link-presence/loss classifier: both a link drop and
a link appearance are treated as evidence, and the fused decision types the
entity crossing between the two nodes --- realising barrier-coverage
geometry on commodity Bluetooth without synchronised nodes.

\textbf{H. Mesh agreement and gamification.} Nearby Ambient Sense instances
exchange crossing records. A crossing is published only once confirmed by a
threshold number of nodes, and rays that disagree with the majority are
down-weighted; the mesh mean becomes the fused estimate, and a reliability
bonus is granted to devices whose measurements lie nearest that mean. A
gamification layer computes an explicit estimated accuracy, rewards
participation (points per crossing), and awards a fusion multiplier when a
crossing is confirmed by several nodes. Because both the reward and the
estimate depend on the mesh mean, rewarding the device nearest that mean is,
by construction, rewarding agreement with the group --- which is a
reasonable proxy for reliability under benign conditions, but is not by
itself resistant to a coordinated group of colluding or spoofed devices;
Section~IV discusses this as a specific direction for hardening. Fig.~\ref{fig:arch}
shows the layered architecture.

\begin{figure}[htb!]
  \centering
  \includegraphics[width=8.7cm]{figures/fig_modern_02_layers.jpg}
  \caption{Layered system architecture of Ambient Sense, from unmodified
  radio hardware through per-link DSP and inference to the application and
  gamification layer.}
  \label{fig:arch}
\end{figure}

\textbf{Privacy and data-bearing boundary.} No camera or microphone is used;
beacon identifiers (MAC, SSID) are not persisted; per-link dBm tables remain
on device; only aggregate crossing records leave the handset, and only when
a mesh exists. The observer is device-free throughout.

\section{Deployment}

Figs.~\ref{fig:sense}--\ref{fig:club} show the deployed
build. The \emph{Sense} tab (Fig.~\ref{fig:sense}) is the live instrument,
reporting passive links across the three bands, people/vehicle/disturbance
counters, a live feature readout, and the attenuation envelope with markers
at recent crossings. The \emph{Scans} tab (Fig.~\ref{fig:scans}) shows a
stored epoch --- window, place, link table, number of links that moved
$\geq 3$~dB, and the resulting field verdict. The \emph{Mesh}
(Fig.~\ref{fig:mesh}), \emph{Transit} (Fig.~\ref{fig:transit}) and
\emph{Club} (Fig.~\ref{fig:club}) tabs expose mesh agreement, the two-node
barrier counter, and gamification scoring, respectively.

\begin{figure}[htb!]
  \centering
  \includegraphics[width=8.7cm]{figures/fig_modern_03_sense.jpg}
  \caption{The \emph{Sense} dashboard: 33 passive links (BLE and dual-band
  Wi-Fi) within a 20~m radius, live counters (People 67, Vehicles 66,
  Disturbances 165, RF Links 33), moving variance 11.45~dB$^2$, jitter
  1.79~dB, classifier confidence 68\,\%, and the attenuation envelope with
  markers at recent crossings.}
  \label{fig:sense}
\end{figure}

\begin{figure}[t]
  \centering
  \includegraphics[width=8.7cm]{figures/fig_modern_04_scans.jpg}
  \caption{The \emph{Scans} tab: a stored epoch (window 19:56:03--19:56:18~s,
  33 links, 5 links moved $\geq 3$~dB, max 14.0~dB) and the verdict ``field
  CHANGED since previous scan'' plus scan history.}
  \label{fig:scans}
\end{figure}

\begin{figure}[t]
  \centering
  \includegraphics[width=8.7cm]{figures/fig_modern_05_mesh.jpg}
  \caption{The \emph{Mesh} tab: 2--5 cooperating nodes confirm each other's
  crossings; the effective radius extends from 4--12~m (single phone) toward
  20--40~m, with a reliability bonus for devices nearest the fused estimate.}
  \label{fig:mesh}
\end{figure}

\begin{figure}[t]
  \centering
  \includegraphics[width=8.7cm]{figures/fig_modern_06_transit.jpg}
  \caption{The \emph{Transit} tab: a two-node barrier over Classic Bluetooth
  RFCOMM and BLE L2CAP; both a link drop and a link appearance are treated as
  evidence, and detection is reported at each separation.}
  \label{fig:transit}
\end{figure}

\begin{figure}[t]
  \centering
  \includegraphics[width=8.7cm]{figures/fig_modern_07_club.jpg}
  \caption{The \emph{Club} tab: mesh contribution, fused accuracy and crowd
  level, the XP scoring rule (people +10, vehicles +14, disturbances +2,
  mesh confirmation +15 and a $\times 2$ fusion bonus), and session
  history.}
  \label{fig:club}
\end{figure}

%---------------------------------------------------------------------------
% III. EVALUATION
%---------------------------------------------------------------------------
\section{Evaluation}

Seven quantities are measured against injected ground truth: (A) per-crossing
sensing accuracy and detection versus distance and radius preset; (B)
scan-epoch movement detection and false-positive rate; (C) ranging coverage
and median absolute error for gated versus ungated estimators; (D) mesh
confirmation and false-confirmation rates; (E) transit-counter detection and
typing versus phone separation; (F) per-band and fused typing accuracy; and
(G) displayed gamified accuracy versus realized accuracy.

All metrics are extracted from a unified, model-based Monte Carlo harness
(1,512 operational scenes and 160 empty control scenes for the main sensing
and mesh experiments, with dedicated sub-experiments of 288 scenes for
ranging and 56 control epochs for change detection, detailed below). The
harness synthesizes radioscape dynamics by combining log-distance path loss,
log-normal scintillation, and Gaussian-profile shadow fading governed by
Fresnel-efficiency decay:
\begin{equation}
P_{\mathrm{rx}}(d) = P_0 - 10n \log_{10}\left(\frac{d}{d_0}\right) - A_0 \exp\left(-\frac{\Delta x^2}{2\sigma^2}\right) \eta_{\mathrm{Fresnel}} + X_\sigma
\label{eq:pathloss}
\end{equation}
Candidate crossings inject explicit entity parameters (pedestrians with
$\sigma = 0.35$~m; vehicles with $\sigma = 1.15$~m, occlusion capped at
34~dB, 5.5~Hz flutter), trajectory vectors, and closest-approach distances
(0.5--20~m) across near (10~m), far (55~m), and mixed AP baselines (33
typical links spanning 12--68~m legs).

Because ground truth is known exactly, accuracy, detection, error, and
false-positive rates are directly measurable, and empty/control scenes are
generated in parallel so that false-change and false-confirmation rates are
measured with denominators rather than assumed to be low. Every candidate
change to a preset, gate, weight vector, epoch rule, or detector fix was A/B
tested on complete scenes before being accepted, and the same distributions
drive what the shipped app shows as radius presets, range bands, and
accuracy figures.

\subsection{Sensing Accuracy and Multi-Band Fusion}
Single-device BLE detection yields 100\,\% sensitivity at 2~m, drops to
$\approx$50\,\% at 14--16~m, and falls below 15\,\% beyond 20~m; the
auxiliary Wi-Fi bands extend physical leg coverage to 58--68~m
(Table~\ref{tab:bands}).

\begin{table}[htb!]
  \centering
  \caption{Per-Band Physical Constants and Measured Per-Crossing Accuracy.}
  \label{tab:bands}
  \footnotesize
  \begin{tabular}{@{}p{1.8cm}p{0.6cm}p{1cm}p{1cm}p{1.2cm}p{1cm}@{}}
    \toprule
    Carrier Band & $n$ & $R_{1\,\mathrm{m}}$ & $A_0$ & Per-Crossing Accuracy & Max Leg \\
    \midrule
    BLE ($2.4$\,GHz)   & $2.4$ & $-42$\,dBm & $7.2$\,dB  & \textbf{57.7\,\%} & $24$\,m \\
    Wi-Fi $2.4$\,GHz   & $2.6$ & $-38$\,dBm & $9.1$\,dB  & $16.1\,\%$        & $58$\,m \\
    Wi-Fi $5$\,GHz     & $2.9$ & $-35$\,dBm & $11.5$\,dB & $25.7\,\%$        & $68$\,m \\
    \bottomrule
  \end{tabular}
\end{table}

Table~\ref{tab:bands} reports each band's \emph{standalone} typing accuracy
when it alone must classify the crossing; this is a different, harder
quantity than the fused, whole-system accuracy reported below, since a
standalone band has no access to the other bands' evidence. BLE outperforms
the Wi-Fi bands standalone (57.7\,\% vs. 25.7\,\% and 16.1\,\%) because it
samples closer to the body and more frequently, even though its 24~m maximum
leg is the shortest of the three. Linearly averaging all three bands instead
of fusing them asymmetrically \emph{degrades} BLE-driven accuracy to
31.0\,\%, which is why the shipped classifier weights BLE for typing and
uses the Wi-Fi bands primarily to extend spatial reach (Section~III.D).

\begin{figure}[htb!]
  \centering
  \includegraphics[width=0.87\linewidth]{Accuracy-Detection.png}
  \caption{BLE per-crossing accuracy/detection (\%) versus crossing distance
  and radius preset (occupied scenes; $N = 13$--$14$ per cell).}
  \label{fig:accuracy}
\end{figure}

\begin{figure}[htb!]
  \centering
  \includegraphics[width=0.87\linewidth]{Eval2.png}
  \caption{Human crossing detection probability and signal attenuation
  across distance and radio band.}
  \label{fig:eval2}
\end{figure}

Per-crossing accuracy across radius presets is shown in
Fig.~\ref{fig:accuracy}. Narrow presets maximise near-field trigger
sensitivity at the expense of field truncation; expansive presets raise the
adaptive trigger threshold to absorb path scintillation over longer
baselines, which blunts sensitivity. The 10~m preset is the recommended
default for standalone operation. At the \emph{fused, whole-system} level
across all 1,512 Monte Carlo runs, pedestrians are correctly typed in
34.9\,\% of occupied scenes and vehicles in 29.6\,\%, while unperturbed
control scenes remain correctly quiescent in 87.3\,\% of cases. Observed
median attenuation dips (Fig.~\ref{fig:eval2}) validate the underlying
physics: shadowing decays with closest-approach distance via Fresnel
efficiency until it intersects the adaptive threshold floor. Reporting
accuracy across discrete radii avoids the misleading bias of global
aggregation: BLE entity typing decays from 92\,\% at 2~m to 47\,\% at 20~m
and 6\,\% at 68~m. Beyond 20--40~m, link-axial Fisher information effectively
collapses~\cite{li2024}; the system therefore reports spatial distance bands
rather than point positions past this range, in line with the physical
channel limit rather than a software choice. In summary, single-device
sensing reliably identifies crossing entities to $\approx$12~m and detects
presence to $\approx$16~m; beyond that, the extended-range Wi-Fi legs serve
as peripheral tripwires rather than sources of fine-grained typing.

\subsection{Scan Epochs and Ranging Estimator Gating}
Inter-epoch analysis tracks RSSI between consecutive observation windows.
Under baseline conditions, movement is detected in 100\,\%, 88\,\%, 81\,\%,
75\,\%, 62\,\%, 12\,\%, 44\,\%, and 25\,\% of candidate crossings at true
distances of 2, 4, 6, 8, 12, 16, 20, and 22~m respectively
(Table~\ref{tab:epoch}), while the average number of perturbed links falls
from 5.0 at 2~m to 0.2 at 22~m.

\begin{table}[htb!]
  \centering
  \caption{Movement Detection and False-Positive Rates for the Scan Epoch Rule.}
  \label{tab:epoch}
  \footnotesize
  \begin{tabular}{@{}lccc@{}}
    \toprule
    Distance & Movement Detected & Naive FPR & Shipped Rule FPR \\
    \midrule
    $2$\,m  & $100\,\%$ & $25\,\%$ & $\mathbf{0.0\,\%}$ \\
    $4$\,m  & $88\,\%$  & $34\,\%$ & $\mathbf{0.0\,\%}$ \\
    $6$\,m  & $81\,\%$  & $41\,\%$ & $\mathbf{0.0\,\%}$ \\
    $8$\,m  & $75\,\%$  & $48\,\%$ & $\mathbf{0.0\,\%}$ \\
    $12$\,m & $62\,\%$  & $55\,\%$ & $\mathbf{0.0\,\%}$ \\
    $16$\,m & $12\,\%$  & $58\,\%$ & $\mathbf{0.0\,\%}$ \\
    $20$\,m & $44\,\%$  & $62\,\%$ & $\mathbf{0.0\,\%}$ \\
    $22$\,m & $25\,\%$  & $65\,\%$ & $\mathbf{0.0\,\%}$ \\
    \bottomrule
  \end{tabular}
\end{table}

The stable-level decision rule (Section~II.F) is what makes this possible: a
naive sample-delta prototype produced 25--62\,\% false-positive alerts under
quiet field conditions, while the shipped 20-sample median comparison drove
false positives to exactly 0.0\,\% across all evaluated distances --- zero
false triggers over 56 control epochs (Rule-of-Three 95\,\% confidence upper
bound: 5.36\,\%).

\begin{table}[htb!]
  \centering
  \caption{Range Estimator (BLE, Per Crossing): Gated ($\ge -75$\,dBm) vs. Ungated.}
  \label{tab:range_estimator}
  \footnotesize
  \begin{tabular}{@{}p{1.4cm}p{1cm}p{1cm}p{1cm}p{1cm}p{1cm}@{}}
    \toprule
    \textbf{AP Layout} & \textbf{Distance} & \textbf{Gated Cov.} & \textbf{Gated Med. $|\mathrm{err}|$} & \textbf{Ungated Cov.} & \textbf{Ungated Med. $|\mathrm{err}|$} \\
    \midrule
    $10$\,m (Near) & $2$\,m  & $100\,\%$ & $2.74$\,m  & $100\,\%$ & $2.74$\,m  \\
    $10$\,m (Near) & $6$\,m  & $83\,\%$  & $1.26$\,m  & $83\,\%$  & $1.26$\,m  \\
    $10$\,m (Near) & $12$\,m & $67\,\%$  & $7.17$\,m  & $67\,\%$  & $7.17$\,m  \\
    $10$\,m (Near) & $22$\,m & $50\,\%$  & $17.19$\,m & $50\,\%$  & $17.19$\,m \\
    $55$\,m (Far)  & $2$--$4$\,m & $\mathbf{0\,\%}$ & \textbf{Refused} & $100\,\%$ & $14$--$23$\,m \\
    \bottomrule
  \end{tabular}
\end{table}

Table~\ref{tab:range_estimator} evaluates 288 scenes across near (10~m) and
far (55~m) AP layouts. On the near layout, every link stays above the
$-75$~dBm gate, so gated and ungated estimates coincide, and median absolute
error grows from 2.74~m at 2~m true distance to 17.19~m at 22~m as the link
weakens. On the far layout, links sit close to the noise floor from the
outset: an ungated estimator reports 100\,\% coverage but with a 14--23~m
error for true crossings at only 2--4~m --- a confident-looking number that
is simply wrong. The $-75$~dBm gate, combined with the valid attenuation
interval $f \in [0.15, 0.92]$, drops far-link coverage to 0\,\%, i.e., the
system refuses to report a point on these links rather than reporting one
that is off by an order of magnitude.

\subsection{Cooperative Mesh and Two-Node Transit Counting}
Confirmation rates in the mesh scale with receiver density, most visibly at
long range (Table~\ref{tab:mesh}, $N=10$ per cell). At 40~m, confirmation
rises from 20\,\% with a single device to 50\,\% across five cooperating
nodes; at 5~m, from 20\,\% to 100\,\%. The non-monotonic step at 10~m
(70\,\%~$\to$~40\,\%) reflects the overhead of multi-node consensus, where
candidate events must align within $\pm5$~s in matching classification
buckets before a crossing is confirmed. Despite this overhead, the voting
logic produced zero false confirmations across all 160 evaluated empty
scenes (Rule-of-Three 95\,\% upper bound: 1.88\,\%), and distributed
coherence extends effective detection from 4--12~m (single phone) toward
20--40~m with mesh cooperation.

\begin{table}[htb!]
  \centering
  \caption{Mesh Confirmation Rate vs. Cooperating Node Count.}
  \label{tab:mesh}
  \footnotesize
  \begin{tabular}{@{}lccc@{}}
    \toprule
    Radius & 1 Node & Intermediate & 5 Nodes \\
    \midrule
    $5$\,m  & $20\,\%$ & --                     & $100\,\%$ \\
    $10$\,m & $70\,\%$ & $40\,\%$ (non-monotonic) & --        \\
    $40$\,m & $20\,\%$ & --                     & $50\,\%$  \\
    \bottomrule
  \end{tabular}
\end{table}

For cross-boundary monitoring, the two-node transit counter (Fig.~\ref{fig:transit})
uses dual Bluetooth Classic RFCOMM and BLE L2CAP links. Evaluated
independently, the BLE channel detects and types 100\,\% of crossings at
3--10~m separations (94\,\% at 15~m), while the RFCOMM loss mechanism alone
delivers 50--56\,\%. The hybrid classifier, which fuses both, sustains
100\,\% detection and typing across separations up to 15~m, decaying
gracefully to 61\,\% at 25~m.

\subsection{Gamified Accuracy and Parameter Selection}
The gamification layer computes a real-time confidence score and structures
rewards as an explicit XP progression: base crossing XP, a $2\times$
multi-band fusion bonus, $+15$~XP for mesh-confirmed events, and a
$1.25\times$ multiplier for nodes closest to the mesh consensus centroid.
Because this reward schedule mirrors the operational likelihood
distribution rather than an arbitrary points table, the displayed accuracy
estimate tracks the same distributions used for the accuracy figures above.

All shipped parameters were selected through the complete-scene A/B protocol
described in Section~I.C, rather than by row-level optimization, which
proved actively misleading: automated weight optimizers marginally improved
row-level $F_1$ (0.173~$\to$~0.179) but degraded under full-scene
simulation --- an unconstrained optimizer converged to a degenerate state
(0\,\% human recall), and a constrained alternative reached only 45.6\,\%
scene accuracy against the shipped heuristic's 46.7\,\%. The same protocol
rejected uniform-weight multi-band averaging (which diluted BLE accuracy to
31.0\,\%, as noted in Section~III.A), noise-floor range extrapolation
(14--23~m phantom offsets, Section~III.B), and naive sample-delta epoch
comparison (25--62\,\% false-positive rates, Section~III.B).

\subsection{Limitations and Untested Conditions}
All results above are model-based, within a simulated test environment;
absolute metrics require empirical on-device re-validation, and real-world
accuracy is expected to degrade somewhat as richer physical multipath is
introduced. Within the modelled envelope, the system identifies crossings to
$\approx$12~m and detects presence to $\approx$16~m, resolving into position
bands rather than discrete points beyond $\approx$8~m; over that 8~m span,
epoch-difference detection operates at zero false positives, and the
mesh-extended configuration reaches a 20--40~m radius at zero false
confirmations.

Several conditions remain untested and define the near-term experimental
agenda on the same harness, pending on-device validation and per-environment
refitting of $n$, $R_{1\text{m}}$, $A_0$, and $\sigma$: indoor settings
(residential 2--5-link layouts, office floors, corridor and stairwell NLOS,
RF-hostile environments, high-density transit hubs with simultaneous
crossings), outdoor settings (urban street canyons, adverse weather, sparse
30--50~m anchor layouts, high-multipath parking lots, rural roads, saturated
events), and practical constraints such as device placement (hand, pocket,
bag), hardware heterogeneity, battery-limited duty cycling, and co-channel
interference. Resilience of the mesh against adversarial beacons and spoofed
reports (Section~I.E) is likewise untested and is treated as future work
below rather than as a delivered capability.

%---------------------------------------------------------------------------
% IV. CONCLUSIONS
%---------------------------------------------------------------------------
\section{Conclusions and Future Work}
Ambient Sense shows that a commodity phone can perform multi-band passive
sensing end to end: fuse BLE and dual-band Wi-Fi features into entity
typing, convert measured dBm into a stated range with an explicit refusal
gate, detect change between two instants of the field without inventing
movement, count transit across a two-node Bluetooth barrier, and extend
accuracy by agreeing with cooperating neighbours --- all validated by a
falsification protocol that A/B tests every shipped constant on complete
scenes rather than row-level metrics. Just as importantly, the evaluation
reports where the underlying physics withdraws its support: 6\,\% typing
accuracy at 68~m, 94\,\% rather than 100\,\% detection at 15~m, and a noise
floor below which ranging is refused rather than merely uncertain.

Because the sensing hardware is the handset the user already carries, the
marginal cost of a new measurement point is near zero, and because the
observer carries no camera, the privacy argument is structural rather than
policy-based. This supports use cases such as ambient-assisted living
(presence and crossing counts without cameras or worn devices), occupancy
and flow analytics in workplaces and corridors, roadside traffic studies via
the two-node barrier counter, and coverage scouting through the
change-detection epochs.

Two directions define the immediate next steps. First, \textbf{on-device
field validation}: every number in Section~III is drawn from a model-based
simulation, and the stated next step is to re-run the same evaluation
protocol on physical hardware across the untested indoor and outdoor
conditions listed in Section~III.E. Second, \textbf{mesh integrity}: the
current reliability bonus rewards agreement with the mesh mean, which is a
reasonable proxy for honesty under benign participation but is not yet a
defence against a coordinated or spoofed minority of devices. Extending it
into a reputation-aware, quality-weighted incentive scheme --- drawing on
established mobile-crowdsensing incentive design~\cite{zhang2022,wang2018}
and informed by known BLE spoofing and downgrade
vulnerabilities~\cite{casar2022,antonioli2020,cai2024} --- is the paper's
main identified direction for future work, alongside broadening the
gamification layer from a point-and-badge system toward published, auditable
reliability scores.

\begin{thebibliography}{21}

\bibitem{wilson2010} J.~Wilson and N.~Patwari, ``Radio tomographic imaging with
wireless networks,'' \emph{IEEE Trans.\ Mobile Comput.}, vol.~9, no.~5,
pp.~621--632, May 2010, doi: 10.1109/TMC.2010.8.

\bibitem{januszkiewicz2018} {\L}.\,Januszkiewicz, ``Analysis of human body
shadowing effect on wireless sensor networks operating in the 2.4~GHz band,''
\emph{Sensors}, vol.~18, no.~10, art.~3412, Oct. 2018, doi:
10.3390/s18103412.

\bibitem{alippi2016} C.~Alippi, M.~Bocca, G.~Boracchi, N.~Patwari, and
M.~Roveri, ``RTI goes wild: Radio tomographic imaging for outdoor people
detection and localization,'' \emph{IEEE Trans.\ Mobile Comput.}, vol.~15,
no.~10, pp.~2585--2598, Oct. 2016, doi: 10.1109/TMC.2015.2504964.

\bibitem{li2024} Z.~Li, A.~Dubey, S.~Shen, N.~K.~Kundu, J.~Rao, and R.~Murch,
``Radio tomographic imaging with reconfigurable intelligent surfaces for
integrated sensing and communication,'' \emph{IEEE Trans.\ Wireless Commun.},
vol.~23, no.~11, pp.~15784--15797, Nov. 2024, doi: 10.1109/TWC.2024.3433011.

\bibitem{wu2024} H.~Wu, C.~Cheng, T.~Peng, H.~Zhou, and T.~Chen, ``Combining
transformer with a latent variable model for radio tomography based robust
device-free localization,'' \emph{Comput.\ Commun.}, vol.~224, art.~108022,
Oct. 2024, doi: 10.1016/j.comcom.2024.108022.

\bibitem{cnnrti2025} ``Passive localization based on radio tomography images
with CNN model utilizing WiFi RSSI,'' \emph{Sci.\ Rep.}, vol.~15, art.~15698,
2025, doi: 10.1038/s41598-025-99694-2.

\bibitem{wong2024} W.~E.~Wong, A.~H.~Wong, W.~Q.~Peh, and C.~K.~Tan, ``A
high-dimensional, multi-transceiver channel state information dataset for
enhanced human activity recognition,'' \emph{Data Brief}, vol.~55,
art.~110673, Aug. 2024, doi: 10.1016/j.dib.2024.110673.

\bibitem{sanchez2020} D.~S\'anchez-Rodr\'iguez, M.~A.~Quintana-Su\'arez,
I.~Alonso-Gonz\'alez, C.~Ley-Bosch, and J.~J.~S\'anchez-Medina, ``Fusion of
channel state information and received signal strength for indoor localization
using a single access point,'' \emph{Remote Sens.}, vol.~12, no.~12,
art.~1995, Jun. 2020, doi: 10.3390/rs12121995.

\bibitem{ramirez2021} R.~Ramirez, C.-Y.~Huang, C.-A.~Liao, P.-T.~Lin, H.-W.~Lin,
and S.-H.~Liang, ``A practice of BLE RSSI measurement for indoor positioning,''
\emph{Sensors}, vol.~21, no.~15, art.~5181, Jul. 2021, doi: 10.3390/s21155181.

\bibitem{milano2024} F.~Milano, H.~da Rocha, M.~Laracca, L.~Ferrigno,
A.~Esp\'iritu Santo, J.~Salvado, and V.~Paciello, ``BLE-based indoor
localization: Analysis of some solutions for performance improvement,''
\emph{Sensors}, vol.~24, no.~2, art.~376, Jan. 2024, doi: 10.3390/s24020376.

\bibitem{bai2020} L.~Bai, F.~Ciravegna, R.~Bond, and M.~Mulvenna, ``A low cost
indoor positioning system using Bluetooth Low Energy,'' \emph{IEEE Access},
vol.~8, pp.~136858--136871, 2020, doi: 10.1109/ACCESS.2020.3012342.

\bibitem{khan2023} D.~Khan and I.~W.-H.~Ho, ``CrossCount: Efficient device-free
crowd counting by leveraging transfer learning,'' \emph{IEEE Internet Things
J.}, vol.~10, no.~5, pp.~4049--4058, Mar. 2023.

\bibitem{zhou2020} R.~Zhou, Z.~Gong, X.~Lu, and Y.~Fu, ``WiFlowCount:
Device-free people flow counting by exploiting Doppler effect in commodity
WiFi,'' \emph{IEEE Syst.~J.}, vol.~14, no.~4, pp.~4919--4930, Dec. 2020.

\bibitem{dmip2023} ``Non-intrusive people counting and identification
simultaneously with commodity WiFi devices,'' in \emph{Proc.\ 6th Int.\ Conf.\
Digital Medicine and Image Processing (DMIP '23)}, 2023, pp.~62--68, doi:
10.1145/3637684.3637694.

\bibitem{zhang2022} J.~Zhang, X.~Li, Z.~Shi, \emph{et al.}, ``A reputation-based
and privacy-preserving incentive scheme for mobile crowd sensing: A deep
reinforcement learning approach,'' \emph{Wirel.\ Netw.}, 2023, doi:
10.1007/s11276-022-03111-9.

\bibitem{wang2018} J.~Wang, M.~Li, Y.~He, H.~Li, K.~Xiao, and C.~Wang, ``A
blockchain based privacy-preserving incentive mechanism in crowdsensing
applications,'' \emph{IEEE Access}, vol.~6, pp.~17545--17556, 2018, doi:
10.1109/ACCESS.2018.2815365.

\bibitem{casar2022} M.~C\"asar, T.~Pawelke, J.~Steffan, and G.~Terhorst, ``A
survey on Bluetooth Low Energy security and privacy,'' \emph{Comput.\ Netw.},
vol.~202, art.~108712, Jan. 2022, doi: 10.1016/j.comnet.2021.108712.

\bibitem{antonioli2020} D.~Antonioli, N.~O.~Tippenhauer, and K.~B.~Rasmussen,
``Key negotiation downgrade attacks on Bluetooth and Bluetooth Low Energy,''
\emph{ACM Trans.\ Priv.\ Secur.}, vol.~23, no.~3, art.~14, Jul. 2020, doi:
10.1145/3394497.

\bibitem{cai2024} H.~Cai, ``Securing billion Bluetooth devices leveraging
learning-based techniques,'' in \emph{Proc.\ AAAI Conf.\ Artif.\ Intell.\
(AAAI)}, vol.~38, no.~21, pp.~23731--23732, 2024, doi:
10.1609/aaai.v38i21.30544.

\bibitem{mobisys2024} ``Poster: Hybrid detection mechanism for spoofing
attacks in Bluetooth Low Energy networks,'' in \emph{Proc.\ 22nd Annu.\ Int.\
Conf.\ Mobile Syst., Appl., Services (MobiSys '24)}, 2024, pp.~710--711, doi:
10.1145/3643832.3661434.

\bibitem{hensler2019} C.~Hensler and P.~Tague, ``DEMO: Using Bluetooth Low Energy
spoofing to dispute device details,'' in \emph{Proc.\ ACM Conf.\ Security
Privacy Wireless Mobile Netw.\ (WiSec '19)}, Miami, FL, USA, 2019, pp.~341--343,
doi: 10.1145/3317549.3326321.

\end{thebibliography}

\end{document}


\end{document}

